\documentclass[letterpaper, 10 pt, conference]{ieeeconf}
\usepackage{amsmath}
\usepackage{amssymb}
\usepackage{graphicx}
\usepackage{booktabs}

\title{ARES: A Lightweight Research Framework for Stable Hybrid Autonomous Navigation}
\author{George David Tsitlauri\\
\textit{Dept. of Informatics \& Telecommunications}\\
\textit{University of Thessaly, Greece}\\
gdtsitlauri@gmail.com}
\date{2026}

\begin{document}
\maketitle

\begin{abstract}
ARES unifies classical control, path planning, distributed coordination, and a
hybrid MPC and RL navigation stack in a lightweight research framework. This
manuscript accompanies the open-source implementation and documents a
reproducible baseline spanning Python, Go, C++, and x86-64 assembly. The
strongest empirical evidence in the current release comes from bounded
cross-language robotics artifacts: sub-microsecond control-kernel benchmarks, a
hybrid navigator that improves on standalone MPC and RL baselines in the
committed scenarios, and fleet-coordination measurements up to 50 robots.
\end{abstract}

\section{Introduction}
ARES targets a gap between heavyweight robotics middleware and small
single-purpose research codebases by providing one compact framework spanning
control, planning, hybrid autonomy, and systems programming.

The design goal is not to replace mature middleware stacks, but to provide a
compact and reproducible experimental substrate where control theory, planning,
systems programming, and learning-based autonomy can be studied together.

\section{Classical Control Theory}
The control module includes PID families, state-space feedback, LQR synthesis,
and Lyapunov-based stability checks. In the current implementation, the
controller family includes standard PID, cascade PID, adaptive gain updates,
fractional-order PID, LQR, H$_\infty$-style robust feedback, $\mu$-synthesis
surrogates, sliding mode control, and backstepping-inspired nonlinear
stabilization. Stability evidence is reported through quadratic certificates,
finite-step decay measurements, and sampled region-of-attraction estimates.

\section{Path Planning Algorithms}
ARES benchmarks graph-search and sampling-based planners on structured grid
maps. The planning suite covers Dijkstra, A*, a compact D* Lite replanning
wrapper, RRT, RRT*, PRM, minimum-snap interpolation, Reeds-Shepp and clothoid
primitives, CHOMP-style smoothing, and TrajOpt-style refinement. A C++ planning
core provides a pybind-ready extension point for accelerated search and
collision checks.

This split between Python orchestration and C++ acceleration is deliberate: the
high-level benchmark logic stays easy to modify, while hot paths can migrate
downward into compiled code as experiments mature.

\section{Real-Time Control in Assembly}
The assembly module implements a PID update kernel in x86-64 NASM using SSE2
floating-point instructions and is paired with a C bridge for benchmarking.
The repository also includes an interrupt simulation routine and a Python
latency benchmark that compares Python, C, and assembly implementations.

\section{Distributed Fleet Coordination}
The Go subsystem models robot state exchange, task assignment, and leader
selection for a lightweight fleet coordinator. The implemented coordinator uses
goroutines for concurrent command generation, an optimal assignment solver based
on the Hungarian algorithm, an RPC service skeleton for robot state exchange,
and a compact Raft-style consensus state for leader election.

In the current environment, the default coordination path remains RPC-based even
though the repository carries a protobuf schema, because the locally available
cached gRPC toolchain expects a newer Go version than the workspace baseline.

\section{ARES-NAVIGATOR Algorithm}
ARES-NAVIGATOR combines a safety-oriented MPC layer with an adaptive learned
policy. The learned policy is trained with a bounded differentiable objective
that remains fast enough for repository-scale verification. A confidence gate
switches between the learned policy and the MPC safety filter, while causal
failure summaries and hybrid Lyapunov certificates support post-hoc analysis.

\section{Experiments}
The repository includes reproducible result artifacts for control response,
planner comparison, assembly latency, fleet coordination, and hybrid navigation.
The committed artifacts include controller response traces, stability summaries,
planner comparisons on $50\times50$, $100\times100$, and $500\times500$ maps,
motion primitive statistics, fleet scaling results, hybrid navigation
comparisons, causal failure counts, and scenario-level simulation outcomes.

\begin{table}[t]
\centering
\caption{Headline Metrics From Committed ARES Artifacts}
\begin{tabular}{ll}
\toprule
Metric & Value \\
\midrule
Assembly PID latency & 180 ns \\
C PID latency & 520 ns \\
Python PID latency & 3800 ns \\
ARES-NAVIGATOR success rate & 0.93 \\
MPC success rate & 0.86 \\
RL success rate & 0.81 \\
Largest fleet benchmarked & 50 robots \\
Fleet overhead at 50 robots & 27.5 ms \\
\bottomrule
\end{tabular}
\end{table}

The current results show the expected qualitative behavior. The sampled
Lyapunov deltas for the LQR and robust-feedback surrogates are negative across
all committed control samples. The assembly path remains the fastest control
kernel in the repository benchmark, and the hybrid ARES-NAVIGATOR policy
outperforms the standalone MPC and RL baselines on the committed success-rate
artifact. On the planning side, the artifact set records graph-search and
sampling-based behavior across three map scales while keeping the experiment
pipeline bounded enough for laptop-class execution.

\subsection{Interpretation Boundary}
The current ARES release should be interpreted as a reproducible robotics
research baseline rather than as a production-ready autonomy stack. Its
strongest claims are local and artifact-driven: the assembly kernel materially
reduces controller latency inside the repository benchmark; the hybrid
navigation layer improves over the current standalone MPC and RL baselines on
the committed scenarios; and the fleet layer supports a bounded 50-robot
evaluation envelope. The planning results remain useful for internal algorithm
comparison, but they are not intended to establish a universal ranking of all
planners across all map regimes.

\section{Limitations and Future Work}
The current repository is intentionally scoped as a reproducible research
baseline rather than a fully productionized robotics stack. Future work should
focus on three areas: first, stronger formal analysis for the hybrid switching
system and robust-control modules; second, a full GPU-scale SAC training path
with longer-horizon environments and replay-based optimization; and third, a
fully enabled gRPC communication layer once the communication stack can be kept
stable within the repository's reproducible environment constraints. Additional
physics fidelity, richer datasets, and hardware-in-the-loop validation would
further strengthen the empirical claims. These remaining items were left open
not because the software baseline could not be extended, but because they
belong to a heavier phase of work involving longer runs, larger systems
integration, and deeper theoretical treatment than the bounded implementation
phase targeted here.

\section{Conclusion}
This implementation establishes the cross-language architecture of ARES and a
reproducible experimentation pipeline. It leaves room for richer simulators,
full gRPC deployment, stronger formal proofs, and hardware-integrated
validation, while already providing a coherent open-source research substrate
with concrete evidence in control latency, hybrid navigation, and bounded fleet
coordination.

\bibliographystyle{IEEEtran}
\end{document}
